\chapter{Ручка азарту}
% full version: chapters/08_nora.tex

Ви вибираєте, де пообідати. Піти в перевірене кафе чи спробувати нове? Якщо завжди ходити в перевірене, ніколи не дізнаєшся, що за рогом є краще. Якщо весь час пробувати нове, часто обідатимеш погано. Де золота середина? Ця задача --- шукати чи користуватися знайденим --- стоїть перед будь-якою істотою, що навчається. У мозку за неї, схоже, відповідає радіостанція \nt{NORA} --- норадреналін.

\section*{Ручка контрастності}

Норадреналін змінює те, наскільки різко нейрон відповідає на сигнал. За слабкої дії нейрон схожий на плавний регулятор яскравості: трохи більший сигнал --- трохи яскравіша відповідь. За сильної він схожий на вимикач: нижче порогу --- тиша, вище --- повна відповідь. Точніше кажучи, досліди показують, що норадреналін пригнічує фонову балаканину нейронів сильніше, ніж їхню відповідь на справжній сигнал; «ручка контрастності» --- зручна модель цього ефекту.

\pencilfig{8}{\pencilart{8}}{Ручка азарту: повернеш ліворуч --- машина пробує нове, праворуч --- рішуче обирає перевірене.}

\section*{Думати чи вирішувати}

Згадайте з третього розділу дві групи, що гальмують одна одну. Коли контрастність низька, обидві працюють, слабша трохи тихіше: мережа ще «думає», зважує варіанти. Коли контрастність вища за певний поріг, мережа різко переходить у режим «вирішила»: одна група перемагає, інша замовкає. Одна ручка перемикає машину між роздумами й рішенням.

\section*{Скільки шукати}

Додамо шум. За низької контрастності шум легко переважує невелику різницю між варіантами, і машина часто пробує не найкращий варіант --- шукає. За високої шум нічого не вирішує, і машина завжди бере те, що вважає найкращим, --- користується.

У нашій моделі світ час від часу змінюється: найкращий варіант стає гіршим. Виграш залежить від ручки як перевернута літера U. Замало контрастності --- машина блукає наосліп. Забагато --- вона не помічає, що світ змінився, і вперто ходить у кафе, яке зіпсувалося. І що частіше змінюється світ, то нижче найкраще налаштування: у мінливому світі вигідніше шукати.

\pencilfig{108}{%
% points: models/nora_explore.py, reward as a share of the best possible, gain 0.5 ... 64 (doubling); the y axis starts at 0.70, hence the break mark
\draw[pen] (0,0) -- (6.3,0); \draw[pen] (0,0) -- (0,0.18); \draw[pen] (0,0.34) -- (0,2.4);
\draw[pcGray, line width=0.6pt] (-0.1,0.14) -- (0.1,0.22); \draw[pcGray, line width=0.6pt] (-0.1,0.3) -- (0.1,0.38);
\node[lbl, anchor=north] at (3.0,-0.05) {ручка контрастності};
\node[lbl, anchor=south, rotate=90] at (-0.1,1.3) {виграш};
\draw[penlite=pcBlue] plot[smooth] coordinates {(0.00,0.55) (0.85,0.79) (1.70,1.19) (2.55,1.68) (3.40,2.00) (4.25,2.07) (5.10,2.01) (5.95,1.91)};
\draw[penlite=pcRed] plot[smooth] coordinates {(0.00,0.55) (0.85,0.69) (1.70,0.93) (2.55,1.24) (3.40,1.40) (4.25,1.28) (5.10,1.06) (5.95,0.89)};
\draw[dotpen=pcBlue, wash=pcBlue] (4.25,2.07) circle (0.07);
\draw[dotpen=pcRed, wash=pcRed] (3.40,1.40) circle (0.07);
\node[lbl, text=pcBlue, anchor=south] at (4.25,2.15) {світ змінюється рідко};
\node[lbl, text=pcRed, anchor=north] at (3.40,1.30) {часто};
\node[lbl, anchor=north west, align=left] at (0.05,-0.35) {блукає\\наосліп}; \node[lbl, anchor=north east, align=right] at (6.3,-0.35) {не помічає\\змін};
}{Виграш моделі за різних положень ручки. Кожна крива має вершину; коли світ змінюється часто, вона нижча й лівіша.}

\section*{Хто крутить ручку}

Природна ознака змін --- несподіванка: помилки стали більшими, ніж зазвичай. За однією гіпотезою, саме про таку несподівану зміну й повідомляє норадреналін. Тут важлива тонкість: сталий, фоновий рівень норадреналіну пов'язаний радше з пошуком і розсіяністю, а короткі спалахи на важливу подію --- із зосередженістю й рішенням. Можливо, спалах працює як переривання в комп'ютері: «кидай поточне, світ змінився, збери план заново».

Непряму ознаку цієї системи видно ззовні: зіниця розширюється, коли щось несподівано змінюється і люди починають учитися швидше. Щоправда, зіниця стежить не лише за норадреналіном, тож це підказка, а не прилад.

І чесне визнання: у нашій моделі підлаштування ручки за несподіванкою виграло зовсім небагато порівняно з найкращим сталим налаштуванням. Де саме таке підлаштування окупається, ще належить з'ясувати.

\fullversion{8}
